\documentclass[10pt]{beamer}

\usepackage{verbatim}
%\usepackage[latin1]{inputenc}
\usepackage{amssymb,amsmath,amsfonts,url}
\usepackage{graphicx,color}
\usepackage{pgf,pgfarrows,pgfnodes,pgfautomata,pgfheaps,pgfshade}
\usepackage[utf8]{inputenc}
\usepackage[french]{babel}

\usepackage[all]{xy}
\usepackage{amstext}
\usepackage{xspace}
\usepackage{mdwtab}


\mode<article> {
  \usepackage{fullpage}
  \usepackage{hyperref}
  \setjobnamebeamerversion{advancedcourse.beamer}
}








\usepackage{listings}
\usepackage{xcolor}

\lstset{
	language=Python,
	basicstyle=\ttfamily\small,
	keywordstyle=\color{blue},
	commentstyle=\color{gray},
	stringstyle=\color{green!50!black},
	showstringspaces=false,
	frame=single,
	numbers=none
}
\lstset{
	language=Python,
	basicstyle=\ttfamily\small,
	keywordstyle=\color{blue},
	commentstyle=\color{gray},
	stringstyle=\color{green!50!black},
	showstringspaces=false,
	frame=single,
	numbers=none,
	literate=
	{é}{{\'e}}1
	{è}{{\`e}}1
	{ê}{{\^e}}1
	{ë}{{\"e}}1
	{à}{{\`a}}1
	{â}{{\^a}}1
	{ä}{{\"a}}1
	{ù}{{\`u}}1
	{û}{{\^u}}1
	{ü}{{\"u}}1
	{ô}{{\^o}}1
	{ö}{{\"o}}1
	{î}{{\^i}}1
	{ï}{{\"i}}1
	{ç}{{\c c}}1
}

















\newcommand{\alertonce}[1]{\addtocounter{beamerpauses}{-1}\alert<+>{#1}}
\newcommand{\alerttwice}[1]{\addtocounter{beamerpauses}{-1}\alert<+-+(1)>{#1}}

%\newcommand{\subsubsubsection}[1]{\noindent\textbf{#1}}

\newcommand{\cfbox}[1]{\begin{center}\fbox{#1}\end{center}}
\newcommand{\cbox}[1]{\begin{center} #1 \end{center}}

%%%%%%%%%%%%%%%%%%%%%%%%%%%%%%%%%%%%%%%%%%%%%%%%%%%%%%%%%%%%%%%%%%%%%%%%%%%%%%%%%%%%%%%%%

\graphicspath{{Graphics/}}
%\input def_0.tex

%\input{StyleSlides.tex}
\mode<presentation> {
    \usepackage{beamerthemeBoadilla}

  \beamertemplatetransparentcovereddynamic
  \beamertemplateballitem
}


\title[IN100 Cours 2]{\textbf{Fondements informatiques I}\\
Cours 2: Structures conditionnelles}

\author[]{Sorina Ionica \texttt{sorina.ionica@uvsq.fr} \\~\\Bilal FAYE  \texttt{bilal.faye@uvsq.fr} }


\institute{}

\date{}


%\setbeamercolor{normal text}{fg=black}

\begin{document}


\frame{\titlepage}



\begin{frame}[containsverbatim]
	\frametitle{Les structures conditionnelles}
	
	\begin{block}{À quoi servent-elles ?}
		Une structure conditionnelle permet d'exécuter certaines instructions
		\textbf{seulement si une condition est vérifiée}.
	\end{block}
	
	\vspace{0.25cm}
	
	
	\begin{exampleblock}{Exemple}
		\begin{center}
			\makebox[0.7\textwidth][l]{%
				\textbf{Si} le solde est suffisant
				$\;\longrightarrow\;$ accepter la transaction.}
			
			\vspace{0.1cm}
			
			\makebox[0.7\textwidth][l]{%
				\textbf{Sinon}
				$\;\longrightarrow\;$ rejeter la transaction.}
		\end{center}
	\end{exampleblock}

	
	\vspace{0.25cm}
	
	\begin{block}{En Python}
		Une condition est une expression qui produit une valeur booléenne :
		\texttt{True} ou \texttt{False}.
	\end{block}
	
	\vspace{0.2cm}
	
	\begin{exampleblock}{Les trois formes que nous allons voir}
		\centering
		\texttt{if}
		\qquad
		\texttt{if ... else}
		\qquad
		\texttt{if ... elif ... else}
	\end{exampleblock}
	
\end{frame}



\begin{frame}[containsverbatim]
	\frametitle{L'instruction \texttt{if}}
	
	\begin{block}{Principe}
		Jusqu'à présent, les instructions étaient exécutées
		\textbf{dans l'ordre}, les unes après les autres.
		
		\vspace{0.1cm}
		
		Avec \texttt{if}, certaines instructions sont exécutées
		\textbf{uniquement si une condition est vraie}.
	\end{block}
	
	\vspace{0.25cm}
	
\begin{alertblock}{Syntaxe}
\begin{lstlisting}[language=Python]
if <condition>:
	instruction1
	instruction2
\end{lstlisting}
\end{alertblock}
	
	\vspace{0.15cm}
	
	\begin{exampleblock}{Bloc d'instructions}
		Les instructions \textbf{indentées} sous le \texttt{if}
		forment un même bloc.
		
		\vspace{0.1cm}
		
		L'indentation permet à Python d'identifier les instructions
		qui appartiennent au bloc conditionnel.
	\end{exampleblock}
	
\end{frame}


\begin{frame}[containsverbatim]
	\frametitle{L'instruction \texttt{if} : un exemple}
	
	
\begin{exampleblock}{Exemple}
\begin{lstlisting}[language=Python]
meteo = "pluie"
if meteo == "pluie":
	print("parapluie")
	print("impermeable")
print("Bonne journée !")
\end{lstlisting}
\end{exampleblock}

	
	\vspace{0.15cm}
	
	\begin{block}{Que se passe-t-il ?}
		\begin{itemize}
			\item \texttt{meteo == "pluie"} vérifie si \texttt{meteo} vaut \texttt{"pluie"}.
			\item Cette condition renvoie \texttt{True} ou \texttt{False}.
			\item Les deux \texttt{print} indentés sont exécutés si la condition est vraie.
			\item \texttt{print("Bonne journée !")} n'est pas indenté : il est exécuté dans tous les cas.
		\end{itemize}
	\end{block}
	
	
\end{frame}


\begin{frame}[containsverbatim]
	\frametitle{L'instruction \texttt{if}-\texttt{else}}
	
	
	\begin{block}{Principe}
		La structure \texttt{if}-\texttt{else} permet de choisir entre
		\textbf{deux blocs d'instructions} :
		
		\begin{itemize}
			\item le bloc \texttt{if} si la condition est vraie ;
			\item le bloc \texttt{else} sinon.
		\end{itemize}
	\end{block}
	
	\vspace{0.2cm}
	
\begin{alertblock}{Syntaxe}
\begin{lstlisting}[language=Python]
if <condition>:
	instruction(s)
else:
	instruction(s)
\end{lstlisting}
\end{alertblock}
	
	
	\vspace{0.15cm}
	
	\begin{exampleblock}{À retenir}
		\begin{itemize}
			\item \texttt{else} est au même niveau d'indentation que \texttt{if}.
			\item \texttt{else} est suivi de deux points \texttt{:}.
			\item Un seul des deux blocs est exécuté.
		\end{itemize}
	\end{exampleblock}

	
\end{frame}
\begin{frame}[containsverbatim]
	\frametitle{L'instruction \texttt{if}-\texttt{else} : un exemple}
	
\begin{exampleblock}{Exemple}
\begin{lstlisting}[language=Python]
meteo = "froid"
if meteo == "pluie":
	print("parapluie")
	print("impermeable")
else:
	print("t-shirt")
	print("shorts")
print("Bonne journée !")
\end{lstlisting}
\end{exampleblock}
	
	
	\vspace{0.15cm}
	
	\begin{block}{Que se passe-t-il ?}
		\begin{itemize}
			\item \texttt{meteo == "pluie"} est faux : le bloc \texttt{else} est exécuté.
			\item Les instructions du bloc \texttt{if} ne sont donc pas exécutées.
			\item \texttt{print("Bonne journée !")} n'est pas indenté :
			il est exécuté après le \texttt{if}-\texttt{else}.
		\end{itemize}
	\end{block}
	
	
\end{frame}
\begin{frame}[containsverbatim]
	\frametitle{L'instruction \texttt{elif}}
	
	
	\begin{block}{Principe}
		\texttt{elif} permet de tester une \textbf{nouvelle condition}
		lorsque la condition précédente est fausse.
		On peut ainsi enchaîner plusieurs conditions.
	\end{block}
	
	
\begin{alertblock}{Syntaxe}
\begin{lstlisting}[language=Python]
if <condition>:
	instruction(s)
elif <condition>:
	instruction(s)
else:
	instruction(s)
\end{lstlisting}
\end{alertblock}
	

	
	\begin{exampleblock}{À retenir}
		\begin{itemize}
			\item \texttt{elif} est au même niveau d'indentation que \texttt{if}. Il est toujours suivi d'une condition et de \texttt{:}. On peut utiliser plusieurs \texttt{elif}. 
			\item \texttt{else}, s'il est présent, est placé à la fin.
		\end{itemize}
	\end{exampleblock}
	
	
\end{frame}


\begin{frame}[containsverbatim]
	\frametitle{L'instruction \texttt{elif} : un exemple}
	
	
\begin{exampleblock}{Exemple}
\begin{lstlisting}[language=Python]
meteo = "froid"
if meteo == "pluie":
	print("parapluie")
	print("impermeable")
elif meteo == "froid":
	print("pull")
	print("echarpe")
else:
	print("t-shirt")
	print("shorts")
print("Bonne journée !")
\end{lstlisting}
\end{exampleblock}
	
\begin{itemize}
	\item Il est possible de mettre autant de \texttt{elif} que l'on souhaite après une condition \texttt{if}.
	\item  Les instructions \texttt{elif} et \texttt{else} sont facultatives : lorsqu'une instruction en \texttt{if} ou \texttt{elif} est définie, il n'est pas obligatoire de prévoir un \texttt{else} après.
\end{itemize}
	
	
\end{frame}

\begin{frame}[containsverbatim]
	\frametitle{Exercice : trouver l'erreur}
	
	
	\begin{block}{Quel est le problème dans ce code ?}
		Observez attentivement l'ordre des instructions et l'indentation.
	\end{block}
	
	\vspace{0.2cm}
	
	\begin{exampleblock}{Code}
		\begin{lstlisting}[language=Python]
meteo = "froid"
if meteo == "froid":
	print("echarpe")
print("Bonne journée !")
elif meteo == "froid" or meteo == "venteux":
	print("manteau")
	print("Bonne journée !")
else:
	print("t-shirt")
print("Bonne journée !")
		\end{lstlisting}
	\end{exampleblock}
	
	
	\begin{alertblock}{Question}
		\centering
		\textbf{Pourquoi ce code provoque-t-il une erreur ?}
	\end{alertblock}
	
	
\end{frame}

\begin{frame}
	\frametitle{Les conditions et les opérateurs de comparaison}
	
	
	\begin{block}{Principe}
		Les opérateurs de comparaison permettent de construire une condition.
		
		Le résultat est toujours \texttt{True} ou \texttt{False}.
	\end{block}
	
	\vspace{0.2cm}
	
	\begin{center}
		\begin{tabular}{|c|l|}
			\hline
			\textbf{Opérateur} & \textbf{Signification} \\
			\hline
			\texttt{<}  & strictement inférieur à \\
			\hline
			\texttt{>}  & strictement supérieur à \\
			\hline
			\texttt{<=} & inférieur ou égal à \\
			\hline
			\texttt{>=} & supérieur ou égal à \\
			\hline
			\texttt{==} & égal à \\
			\hline
			\texttt{!=} & différent de \\
			\hline
		\end{tabular}
	\end{center}
	
	\vspace{0.2cm}
	
	\begin{exampleblock}{Exemple}
		\centering
		\texttt{a >= b}
		$\;\longrightarrow\;$ \texttt{True} si $a$ est supérieur ou égal à $b$,
		\texttt{False} sinon.
	\end{exampleblock}
	
	
\end{frame}

\begin{frame}[containsverbatim]
	\frametitle{Exemple}
	

\begin{exampleblock}{Vérifier si une transaction est possible}
\begin{lstlisting}[language=Python]
solde = 20.0
prix_achat_hors_taxe = 19.0
taxe = 1.0
if solde >= prix_achat_hors_taxe * taxe:
	print("Transaction possible")
else:
	print("Transaction impossible")
print("Fin")
\end{lstlisting}
\end{exampleblock}
	
	
	\vspace{0.1cm}
	
	\begin{block}{À retenir}
		Dans la condition, la multiplication est effectuée avant la comparaison :
		\[
		19.0 \times 1.08 = 20.52
		\]
		Donc $20.0 \geq 20.52$ est faux.
	\end{block}
	
	
\end{frame}
\begin{frame}[containsverbatim]
	\frametitle{Priorité des opérateurs}
	
	
	\begin{center}
		\begin{tabular}{|c|}
			\hline
			\textbf{Du plus prioritaire au moins prioritaire} \\
			\hline
			\texttt{()} \\
			\hline
			\texttt{**} \\
			\hline
			\texttt{*}, \texttt{/}, \texttt{//}, \texttt{\%} \\
			\hline
			\texttt{==}, \texttt{!=}, \texttt{<}, \texttt{>}, \texttt{<=}, \texttt{>=} \\
			\hline
			\texttt{not} \\
			\hline
			\texttt{and} \\
			\hline
			\texttt{or} \\
			\hline
		\end{tabular}
	\end{center}
	
	\vspace{0.15cm}
	
	\begin{exampleblock}{Exercice}
Pour $a=3$ et $b=2$, que valent les expressions suivantes ?

\begin{lstlisting}[language=Python]
a, b = 3, 2
a + b * 2
a > 0 and b * 2 == 4
not a > 0 or b > 5
not (a > 0 and b > 5)
2 ** 3 // 4
17 // 5 + 17 % 5
\end{lstlisting}
\end{exampleblock}
	
\end{frame}
\begin{frame}[containsverbatim]
	\frametitle{Évaluation paresseuse avec \texttt{and} et \texttt{or}}

	\begin{block}{Principe}
		Python évalue les expressions de gauche à droite et
		\textbf{s'arrête dès que le résultat est connu}.
	\end{block}
	
	\vspace{0.2cm}
	
	\begin{columns}[T]
		\begin{column}{0.48\textwidth}
			\begin{alertblock}{\texttt{x and y}}
				Si \texttt{x} est \texttt{False},
				\texttt{y} n'est pas évalué.
				
				\vspace{0.1cm}
				
				\centering
				\texttt{False and ...}
			\end{alertblock}
		\end{column}
		
		\begin{column}{0.48\textwidth}
			\begin{alertblock}{\texttt{x or y}}
				Si \texttt{x} est \texttt{True},
				\texttt{y} n'est pas évalué.
				
				\vspace{0.1cm}
				
				\centering
				\texttt{True or ...}
			\end{alertblock}
		\end{column}
	\end{columns}
	
	\vspace{0.15cm}
	
\begin{exampleblock}{Exemple}
\begin{lstlisting}[language=Python]
False and (10 / 0)   # pas d'erreur
True or (10 / 0)     # pas d'erreur
\end{lstlisting}
\end{exampleblock}
	
\end{frame}
\begin{frame}[containsverbatim]
	\frametitle{Exemples avec \texttt{and}}
	
	
	\begin{exampleblock}{Premier cas}
\begin{lstlisting}[language=Python]	
x = 0
if x != 0 and 2 // x == 2:
	print("Test réussi")
print("Fin")
\end{lstlisting}
		\centering
		\textbf{Pas d'erreur :} \texttt{x != 0} est faux, donc la division n'est pas évaluée.
	\end{exampleblock}
	
	
	\vspace{0.15cm}
	
	\begin{exampleblock}{Deuxième cas}
\begin{lstlisting}[language=Python]
x = 0
if 2 // x == 2 and x != 0:
	print("Test réussi")
print("Fin")
\end{lstlisting}
\centering
\textbf{Erreur :} la division est évaluée en premier.
\end{exampleblock}
	
\end{frame}
\begin{frame}[containsverbatim]
	\frametitle{Exemples avec \texttt{or}}
	
	
	\begin{exampleblock}{Premier cas}
\begin{lstlisting}[language=Python]
x = 0
if x == 0 or 2 // x == 2:
	print("Test réussi")
print("Fin")
\end{lstlisting}
		\centering
		\textbf{Pas d'erreur :} \texttt{x == 0} est vrai,
		donc la division n'est pas évaluée.
	\end{exampleblock}
	
	

	
	\begin{exampleblock}{Deuxième cas}
\begin{lstlisting}[language=Python]
x = 0
if 2 // x == 2 or x == 0:
	print("Test réussi")
\end{lstlisting}
		\centering
		\textbf{Erreur :} la division est évaluée en premier.
	\end{exampleblock}
	
	

	
	\begin{block}{À retenir}
		Avec \texttt{and} comme avec \texttt{or}, l'ordre des conditions
		peut changer le résultat de l'évaluation.
	\end{block}
	
	
\end{frame}
\begin{frame}[containsverbatim]
	\frametitle{La portée des variables}
	
	
	\begin{block}{Définition}
		La \textbf{portée} (\textit{scope}) d'une variable désigne
		les endroits du programme où l'on peut utiliser cette variable.
	\end{block}
	
	\vspace{0.15cm}
	
	\begin{exampleblock}{Exemple}
\begin{lstlisting}[language=Python]
num1 = 1
num2 = 2
resultat = "résultat non connu"
if num1 < num2:
	resultat = "num2 est plus grand que num1"
print(resultat)   # resultat est toujours accessible
\end{lstlisting}
	\end{exampleblock}
	
	
	\vspace{0.1cm}
	
	\begin{alertblock}{À retenir}
		En Python, un bloc \texttt{if} ne crée pas une nouvelle portée :
		\texttt{resultat} reste accessible après le \texttt{if}.
	\end{alertblock}

	
\end{frame}

\begin{frame}[containsverbatim]
	\frametitle{La portée des variables : attention}
	
	
\begin{exampleblock}{Que se passe-t-il ?}
\begin{lstlisting}[language=Python]
num1 = 2
num2 = 1
if num1 < num2:
	maximum = num2
print(maximum)
\end{lstlisting}
\end{exampleblock}
	
	
	\vspace{0.15cm}
	
	\begin{alertblock}{Erreur}
		La condition est fausse : le bloc \texttt{if} n'est pas exécuté.
		
		\texttt{maximum} n'est donc jamais créée.
		
		$\rightarrow$ \texttt{print(maximum)} provoque une erreur.
	\end{alertblock}
	
	
\end{frame}


\begin{frame}[containsverbatim]
	\frametitle{Instructions \texttt{if} imbriquées}
	

	\begin{block}{Principe}
		Un \texttt{if} peut être placé à l'intérieur d'un autre \texttt{if}.
		Le second \texttt{if} est alors exécuté seulement si le premier est vrai.
	\end{block}
	
	\vspace{0.15cm}
	
	\begin{exampleblock}{Exemple}
		\begin{lstlisting}[language=Python]
num1, num2 = 2, 5
maximum = num1
if num1 < num2:
	maximum = num2
	num3 = 3
	if num2 < num3:
		maximum = num3
print(maximum)
		\end{lstlisting}
	\end{exampleblock}
	
	
	
	\begin{alertblock}{À vous de réfléchir}
		Que vaut \texttt{maximum} à la fin ?
		
		À quel endroit \texttt{num3} est-elle créée ?
	\end{alertblock}
	
	
\end{frame}



\begin{frame}[containsverbatim]
	\frametitle{Qu'affiche ce programme ?}

	
\begin{exampleblock}{}
\begin{lstlisting}[language=Python]
solde = 21.0
taxe = 1.08
prix = 19.0
nom_carte = "Daniel Dupont"
nom_client = "Daniel Dupont"
vendeur = "Fred"
vendeurs = ["Maria", "Yoann", "Emilie", "Tia"]
decouvert = True
if solde <= prix * taxe and not decouvert:
	print("Transaction refusée")
else:
	if nom_carte != nom_client:
		print("Client non valide")
	else:
		if vendeur not in vendeurs:
			print("Vendeur non autorisé")
		else:
			print("Transaction approuvée")
\end{lstlisting}
\end{exampleblock}
	
	\centering
	\textit{Suivez les conditions une par une pour trouver le résultat.}
	
	
\end{frame}





\end{document}

\frame[containsverbatim]{\frametitle{Pour aller plus loin ...}

\begin{alertblock}
\begin{verbatim}
meteo="pluie"
match meteo: 
	case "pluie": 
		print("parapluie")
	case "soleil"; 
		print("t-shirt")
	case "froid"
		print("manteau")			
\end{verbatim}
\end{alertblock}


